\documentclass[11pt]{letter}
\usepackage[a4paper, margin=1in]{geometry}
\usepackage{hyperref}
\usepackage{parskip}

\signature{Lizhuo Zhang, Ph.D.\\
Hunan Agricultural University\\
\href{mailto:zhanglizhuo@hunau.edu.cn}{zhanglizhuo@hunau.edu.cn}}
\address{Lizhuo Zhang\\
Hunan Agricultural University\\
Changsha, Hunan, China}

\begin{document}
\begin{letter}{The Editors\\PLOS ONE}

\opening{Dear Editors,}

We are pleased to submit our manuscript, \textit{``A Four-Dimension
Pre-Modeling Audit Protocol for Tabular Prediction Benchmarks: Evidence
from Seven Public Educational Datasets,''} for consideration by PLOS ONE
as an Original Research article.

The manuscript proposes and empirically validates a reusable
\textit{pre-modeling} audit protocol for tabular prediction benchmarks.
The protocol integrates four well-established diagnostics into a single
operational checklist: (i) baseline gap, (ii) split instability,
(iii) permutation-based null separation, and (iv) metadata adequacy
under group-aware holdout. Although each tool is standard in isolation,
their systematic combination---with explicit operational thresholds and
profile labels---and the cross-dataset evidence from seven public
educational releases ($N = 145$ to $32{,}593$) are, to our knowledge, new.
The empirical findings are striking: five of seven widely used educational
datasets fail group-aware holdout under linear models, and three of those
remain fragile even for ensemble models. We believe these results have
direct implications for any tabular benchmark with nested group structure,
not only educational ones.

\textbf{Why PLOS ONE.} The paper's primary value is methodological rigor
on a domain-relevant question rather than a single domain
``breakthrough.'' PLOS ONE's editorial criteria---technical soundness,
methodological clarity, and ethical standards rather than perceived
novelty or impact---match the contribution well. Many of our findings are
\textit{negative} in nature (most evaluated datasets fail at least one
audit dimension); PLOS ONE has a long-standing commitment to publishing
rigorous negative and replication-flavored work.

\textbf{Prior submission history.} We previously submitted an earlier
version of this manuscript to \textit{Education Sciences} (manuscript ID
education-4295342). The submission was returned at the editorial stage as
out-of-scope for the journal's Higher Education section, without
peer review. After receiving that decision we revised the framing to
emphasize the methodological (cross-domain) contribution rather than a
specific education sub-field, expanded the technical disclosure
(see below), and re-targeted the manuscript to a venue whose criteria
align with the contribution. We disclose this history in the interest of
full transparency. The manuscript is not under consideration at any other
journal.

\textbf{Revisions since the prior submission.} (i) The permutation-null
analysis now stores the full empirical $R^2$ distribution (500 draws
across 10 splits per dataset) so that the corresponding figure plots the
real null directly rather than a parametric reference; (ii) the MM-TBA
case now explicitly reports the disagreement between MAE-based beat rate
($0.80$) and iid $R^2$ ($-0.067$) and explains why both signals support
the Fragile profile; (iii) the title and abstract emphasize the
\textit{tabular prediction benchmark} framing, with educational datasets
positioned as the empirical testbed.

\textbf{Data and code availability.} All seven datasets are public.
The audit protocol implementation, raw run outputs (including
permutation-null distributions), and figure-regeneration scripts are
available in a public GitHub repository linked in the Data Availability
statement.

\textbf{Authorship and conflicts.} All listed authors meet ICMJE
authorship criteria. The authors declare no competing interests.
The study used existing public dataset releases only; no new human-subjects
data were collected, and ethical review was waived under our institutional
policy (statement included in the manuscript).

\textbf{Suggested academic editors / reviewers.} Editors with expertise in
educational data mining, learning analytics, or applied
machine-learning methodology would be most appropriate. We are happy to
provide reviewer suggestions on request.

We thank you for considering this work and look forward to your decision.

\closing{Sincerely,}

\end{letter}
\end{document}
